Classical pathfinding assumes a static environment in which the agent can only choose where to move. In many practical settings -- warehouse robots pushing crates aside, excavation planning, or game agents that interact with a destructible world -- the agent can also \emph{modify} the cost landscape it traverses. We refer to this class of problems as \emph{Pathfinding in Interactive Environments} (PIE). PIE is naturally multi-objective: shifting weight away from a high-cost cell shortens the traversal but pays a price in modifications, and the user of the planner generally wants to inspect a Pareto front of step-count vs.\ shift-cost trade-offs rather than a single optimum.

This paper introduces \textbf{MCTS-PIE}, a comparative study of Monte Carlo Tree Search variants for PIE, alongside an $\varepsilon$-constraint A* baseline and an A*+MCTS hybrid. We study six algorithms under a uniform protocol -- 30 seeds across 10 maps -- and report results in terms of standard multi-objective performance indicators. The contributions are: (i) a formalization of PIE as a multi-objective sequential decision problem, (ii) single- and multi-objective MCTS variants with global Pareto archives, (iii) a two-phase MCTS that decomposes path and shift search, (iv) an A*+MCTS hybrid, and (v) a reproducible experimental harness with full per-run Pareto archives, 30-seed coverage, and Holm--Bonferroni-corrected statistics.
